% PyOmniGraph Paper — CEUR-WS submission (skeleton)
% Open source at https://github.com/WolfgangFahl/pyomnigraphPaper
%
% Target venue: 6th Wikidata Workshop 2026 (co-located with ISWC 2026, Bari),
% Novel Work track, published in CEUR-WS. In EasyChair, prefix the submission
% title field with "[Novel]".
%
% CEUR-WS 1-column style (ceurart). Do NOT use the twocolumn / hf options
% for papers submitted to CEUR-WS.
%
% CONTENT POLICY (see AGENTS.md):
%   This is a scientific paper targeting CEUR-WS. AI assistance is limited
%   to spelling/grammar and formatting/build tooling. AI must NOT generate,
%   rewrite or contribute scientific content, arguments or prose.
%   All sections below are intentionally left EMPTY for the authors to write.

\documentclass[
% twocolumn, % NOT allowed for CEUR-WS
% hf,        % NOT allowed for CEUR-WS
]{ceurart}

\sloppy

% Minted/listings support (optional)
\usepackage{listings}
\lstset{breaklines=true}

\begin{document}

%% Rights management information (CC-BY is the CEUR-WS default)
\copyrightyear{2026}
\copyrightclause{Copyright for this paper by its authors.
  Use permitted under Creative Commons License Attribution 4.0
  International (CC BY 4.0).}

%% Conference information
\conference{6th Wikidata Workshop 2026, co-located with ISWC 2026,
  October 25--26, 2026, Bari, Italy}

%% Title
% TODO(authors): finalize the paper title.
\title{PyOmniGraph - Getting and hosting your own federated copies of subgraphs of Wikidata and other knowledge graphs}
\tnotemark[1]
\tnotetext[1]{Preprint. Intended for submission to the 6th Wikidata Workshop 2026
  (co-located with ISWC~2026); if accepted, to appear in the CEUR-WS proceedings.
  This version has not been peer-reviewed.}

%% Authors and affiliations
% Confirmed authors only. ORCIDs taken from the media wiki scholar pages.
% Further co-authors are NOT listed until participation is confirmed.
\author[1]{Wolfgang Fahl}[%
orcid=0000-0002-0821-6995,
email=wf@bitplan.com,
]
\cormark[1]
\author[1]{Tim Holzheim}[%
orcid=0000-0003-2533-6363,
]
\author[2]{Christoph Lange}[%
orcid=0000-0001-9879-3827,
]
\author[1]{Stefan Decker}[%
orcid=0000-0001-6324-7164,
]

\address[1]{RWTH Aachen University, Computer Science i5, Aachen, Germany}
\address[2]{Fraunhofer FIT, Sankt Augustin, Germany}

\cortext[1]{Corresponding author.}

%% Abstract (author-provided text).
\begin{abstract}
Wikidata is an example for a very large, crowd sourced general knowledge graph backed by a world wide community since 2012. Factgrid is using the same Wikibase infrastructure as Wikidata since 2017 for historical research projects. GOV (Geschichtliches/Genealogisches Ortsverzeichnis) is an online gazetteer for searching ancient places. Wikidata grew so large that the Blazegraph endpoint powering it was pushed to the 4 Terrabyte limit and the Wikimedia Foundation therefore decided to split the knowledge graph to host two instances. The challenges of federated queries that came up in the above use cases led to the need of being able to work with subgraphs of the RDF graphs and try out different alternative Endpoints such as QLever and Virtuoso and others. 
While typical research projects that do such subgraphing and benchmarks roll their own environments we decided to provide a python based  Open Source library pyOmnigraph to provide a unified infrastructure with a standard set of REST APIs and CLI commands to improve the handling and potentially get a bigger matrix of subgraphs versus endpoints to work with. 
In this work we report on pyOmnigraph usage on a matrix of three RDF graphs: Wikidata, Factgrid and GOV and endpoints blazegraph, jena, Qlever on a selection of subgraph combinations.

\end{abstract}

%% Keywords
\begin{keywords}
  Wikidata \sep
  FactGrid \sep
  QLever \sep
  Blazegraph \sep
  Apache Jena \sep
  GOV \sep
  RDF triple store \sep
  federated SPARQL \sep
  benchmark \sep
  PyOmniGraph
\end{keywords}

\maketitle

\section{Introduction}\label{sec:introduction}
Wikidata~\cite{Vrandecic2014} started in 2012 as an effort to back the different international Wikipedia sites with a common set of facts. Initially, some 2 million entities have been harvested from existing Wikipedia content. As of 2016, over 100 million entities have been curated by a growing worldwide community of contributors. In 2022, we reported on How to get and host your own copy of Wikidata~\cite{fahl2022wikidata}. 


% From a single Wikidata copy to your own copies of many RDF sources.
% Intentionally empty — to be written by the authors.

\section{Related Work}\label{sec:relatedwork}
% Predecessor 2022 paper, RDF triple stores, SPARQL 1.1 federation, benchmarking.
% Intentionally empty — to be written by the authors.

\section{The PyOmniGraph Resource}\label{sec:resource}
% Unified Python API, Docker deployment, YAML configuration.
% Intentionally empty — to be written by the authors.

\subsection{Supported Backends}\label{sec:backends}
% Jena, Blazegraph, Oxigraph, QLever (working); GraphDB, MillenniumDB,
% Virtuoso, Stardog (planned). Intentionally empty.

\subsection{Command-Line Interface}\label{sec:cli}
% omnigraph (server management), rdfdump (data operations). Intentionally empty.

\section{Data Sources}\label{sec:datasources}
% Subgraphs of Wikidata, FactGrid, GOV, and potentially further RDF sources.
% Intentionally empty — to be written by the authors.

\section{Federated Query Benchmark}\label{sec:benchmark}
% Combinations of endpoints; query set; metrics (load time, query latency and
% throughput); reproducibility. Intentionally empty — to be written by the authors.

\section{Evaluation}\label{sec:evaluation}
% Results across store and endpoint combinations. Intentionally empty.

\section{Availability}\label{sec:availability}
% Apache-2.0, PyPI, GitHub; resource metadata. Intentionally empty.

\section{Conclusion and Future Work}\label{sec:conclusion}
% Intentionally empty — to be written by the authors.

%% Declaration on Generative AI (required by CEUR-WS, see
%% https://ceur-ws.org/GenAI/Policy.html). Finalize per the AGENTS.md
%% governance in this repository before submission.
\section*{Declaration on Generative AI}
AI assistance in this project is limited to spelling/grammar and formatting/build tooling; AI is NOT used to
generate scientific content, arguments or prose.

%% Bibliography
\bibliography{references}

\end{document}
